\documentclass[11pt,a4paper]{article}

\usepackage[margin=1in]{geometry}
\usepackage{amsmath,amssymb,amsfonts}
\usepackage{bm}
\usepackage{booktabs}
\usepackage{algorithm}
\usepackage{algpseudocode}
\usepackage{hyperref}

\title{EVA Channel Simulation: System Model, Channel Estimation,\\
Demodulation, and Maximum Likelihood Estimation}
\author{OTFS/OFDM Research Project}
\date{\today}

\newcommand{\deltaf}{\Delta f}
\newcommand{\CN}{\mathcal{CN}}

\begin{document}
\maketitle

\section{System Parameters}

The simulation targets the Extended Vehicular~A (EVA) channel model at a
carrier frequency $f_c = 4$\,GHz with subcarrier spacing $\deltaf = 15$\,kHz.
The time--frequency (TF) grid has $N = 512$ subcarriers and $M = 128$ symbols,
giving a symbol duration $T = 1/\deltaf$ (no cyclic prefix, $T_s = T$). Data
symbols are drawn from a unit-energy 4-QAM constellation. A mobile
terminal velocity of $120$\,km/h produces a maximum Doppler shift
$\nu_{\max} = v f_c / c \approx 444$\,Hz, corresponding to $k_\nu =
\lceil \nu_{\max} M / \deltaf \rceil = 4$ Doppler bins.

\section{Multipath Channel Model}

The baseband channel is modeled as a sum of $P$ discrete specular paths,
each characterized by a complex gain $\alpha_p$, a delay $\tau_p$ and a
Doppler shift $\nu_p$:
\begin{equation}
  h(\tau,\nu) \;=\; \sum_{p=1}^{P} \alpha_p\,
  \delta(\tau - \tau_p)\,\delta(\nu - \nu_p).
\end{equation}

\subsection{Nine-Path (full EVA) Version}

The standard EVA profile has $P = 9$ paths with the delays and relative
average powers listed in Table~\ref{tab:eva9}. The integer delay tap for
path $p$ is $l_p = \operatorname{round}(\tau_p N \deltaf)$, producing the
sequence $[0,0,1,2,3,5,8,13,19]$. The maximum delay tap is
$l_\tau = \lceil \tau_{\max} N \deltaf \rceil = 20$.

\begin{table}[h]
\centering
\caption{Full EVA channel (9 paths).}
\label{tab:eva9}
\begin{tabular}{rrr}
\toprule
Path & Delay $\tau_p$ [ns] & Rel. Power [dB] \\
\midrule
1 & 0    &  0.0 \\
2 & 30   & $-1.5$ \\
3 & 150  & $-1.4$ \\
4 & 310  & $-3.6$ \\
5 & 370  & $-0.6$ \\
6 & 710  & $-9.1$ \\
7 & 1090 & $-7.0$ \\
8 & 1730 & $-12.0$ \\
9 & 2510 & $-16.9$ \\
\bottomrule
\end{tabular}
\end{table}

\subsection{Three-Path (reduced) Version}

For the reduced simulation, only the three strongest and well-separated
taps are retained, as shown in Table~\ref{tab:eva3}. This variant keeps
the essential multipath diversity while reducing the number of unknown
channel parameters to estimate, which simplifies MLE/CRB analysis and
shortens simulation time.

\begin{table}[h]
\centering
\caption{Reduced EVA channel (3 paths).}
\label{tab:eva3}
\begin{tabular}{rrr}
\toprule
Path & Delay $\tau_p$ [ns] & Rel. Power [dB] \\
\midrule
1 & 0   &  0.0 \\
2 & 150 & $-1.4$ \\
3 & 370 & $-0.6$ \\
\bottomrule
\end{tabular}
\end{table}

In both versions the relative powers are converted to linear scale and
normalized to sum to unity; the path amplitudes used in simulation are
$|\alpha_p| = \sqrt{\bar P_p}$ multiplied by a uniformly random phase
$e^{j\phi_p}$. At every Monte Carlo trial each path's Doppler shift is
redrawn as $\nu_p = \nu_{\max}\cos(\theta_p)$ with
$\theta_p \sim \mathcal{U}(-\pi,\pi)$, giving the classical Jakes law.

\subsection{Time--Frequency Input/Output Relation}

With ICI (inter-carrier interference) retained, the TF-domain channel
applied in \texttt{channel.m} maps $X[n,m]$ to
\begin{equation}
  Y[n,m] \;=\; \sum_{p=1}^{P}\alpha_p
  \sum_{n'=0}^{N-1}
  I_{n-n'}(\nu_p)\;
  X[n',m]\;
  e^{-j2\pi\,\deltaf\,n'\tau_p}\;
  e^{\,j2\pi T_s m\,\nu_p}
  \;+\; Z[n,m],
  \label{eq:channel_tf}
\end{equation}
where $Z[n,m]\sim\CN(0,N_0)$ is AWGN and $I_k(\nu)$ are the ICI
coefficients produced by the Doppler-induced loss of orthogonality
(implemented by \texttt{get\_ici\_coeffs.m}). The Monte-Carlo
simulations always propagate the transmit signal through
Eq.~\eqref{eq:channel_tf} — the full ICI kernel is retained, so the
Doppler-induced loss of subcarrier orthogonality is part of every
received waveform.

When $\nu_p T \ll 1$ the ICI kernel becomes nearly diagonal
($I_0 \approx \mathrm{sinc}(\nu T)$), giving the approximation
\begin{equation}
  Y[n,m] \;\approx\; \sum_{p=1}^{P}\alpha_p\, X[n,m]\,
  e^{-j2\pi \deltaf n \tau_p}\, e^{\,j2\pi T_s m \nu_p} + Z[n,m].
  \label{eq:channel_noici}
\end{equation}
Equation~\eqref{eq:channel_noici} is \emph{not} used to generate data,
and it is \emph{not} used for the MCRB computation either. The MCRB
routine (\texttt{getMCRB\_TF.m}) builds its mean and all gradients/
Hessians with the full ICI kernel $I_k(\nu)$ and its first and second
derivatives (see \texttt{get\_ici\_coeffs\_derivative*.m}). The
no-ICI form in Eq.~\eqref{eq:channel_noici} is only the signal model
assumed \emph{inside the MLE}: it lets the estimator separate each
$(\tau_p,\nu_p)$ pair through a simple rank-one basis, at the price of
a model mismatch with respect to the ICI-bearing observation. This
mismatch is exactly what the MCRB characterizes.

\section{OFDM Transmission}

\subsection{Signal Model and CSI Demodulation}
OFDM places data symbols $X[n,m]$ directly on the TF grid. With perfect
CSI the receiver inverts the full TF channel column by column. Each
OFDM symbol $m$ yields a linear system
\begin{equation}
  \mathbf{y}_m \;=\; \mathbf{H}_m\mathbf{x}_m + \mathbf{z}_m,
\end{equation}
where $\mathbf{H}_m \in \mathbb{C}^{N\times N}$ is assembled from
Eq.~\eqref{eq:channel_tf} and captures the Doppler-induced ICI. The
demodulator applies a single-shot MMSE equalizer
\begin{equation}
  \hat{\mathbf{x}}_m \;=\;
  \bigl(\mathbf{H}_m^{\mathsf H}\mathbf{H}_m + N_0\mathbf{I}\bigr)^{-1}
  \mathbf{H}_m^{\mathsf H}\mathbf{y}_m,
\end{equation}
followed by nearest-neighbor 4-QAM slicing. This is implemented in
\texttt{channel\_inverse.m}.

\subsection{Channel Estimation with Scattered Pilots}
Channel estimation for OFDM uses a scattered pilot sublattice on the TF
grid with spacing $(\Delta n_p, \Delta m_p) = (8,8)$, chosen to
oversample the delay/Doppler spread
($\Delta n_p < N/(2 l_\tau)$, $\Delta m_p < M/(2 k_\nu)$). At every
pilot $(n_i,m_j)$ a least-squares one-tap estimate is formed,
\begin{equation}
  \hat H_{\mathrm{LS}}[n_i,m_j] \;=\;
  \frac{Y[n_i,m_j]}{X[n_i,m_j]},
\end{equation}
and the full TF channel map is reconstructed by 2-D linear
interpolation with nearest-neighbor extrapolation
(\texttt{griddedInterpolant}).

\paragraph{2-D bilinear interpolation.}
The LS estimates are available on the regular sublattice
$\{n_i\}\times\{m_j\}$ with $n_i = 1,1+\Delta n_p,\dots$ and
$m_j = 1,1+\Delta m_p,\dots$. For an arbitrary TF bin $(n,m)$ that
falls between four neighbouring pilot positions
$(n_i,m_j)$,\,$(n_i,m_{j+1})$,\,$(n_{i+1},m_j)$,\,$(n_{i+1},m_{j+1})$,
define the normalized coordinates
\begin{equation}
  t \;=\; \frac{n - n_i}{n_{i+1} - n_i}, \qquad
  u \;=\; \frac{m - m_j}{m_{j+1} - m_j},
  \qquad t,u\in[0,1].
\end{equation}
The bilinear interpolant is
\begin{equation}
  \hat H[n,m] \;=\;
  (1-t)(1-u)\,\hat H_{\mathrm{LS}}[n_i,m_j]
  \;+\; t(1-u)\,\hat H_{\mathrm{LS}}[n_{i+1},m_j]
  \;+\; (1-t)u\,\hat H_{\mathrm{LS}}[n_i,m_{j+1}]
  \;+\; tu\,\hat H_{\mathrm{LS}}[n_{i+1},m_{j+1}].
\end{equation}
At the grid edges, where $(n,m)$ lies outside the pilot sublattice
(i.e.\ $n < n_1$ or $n > n_{N_p}$, and similarly for $m$), the
nearest-neighbor extrapolation strategy clamps the coordinates to the
boundary pilots, effectively copying the closest interpolated value
outward. This avoids artefacts from polynomial extrapolation in regions
where no surrounding pilot rectangle exists.

Data symbols are recovered with a
per-bin 1-tap MMSE equalizer
\begin{equation}
  \hat X[n,m] \;=\;
  \frac{\hat H^{*}[n,m]}{|\hat H[n,m]|^2 + N_0}\,Y[n,m].
\end{equation}
Because this ignores ICI, a BER floor appears at high SNR.

\section{OTFS Transmission}

OTFS places data in the delay--Doppler (DD) grid. The
inverse symplectic finite Fourier transform (\texttt{isfft}) converts
DD to TF before transmission, and the symplectic finite Fourier
transform (\texttt{sfft}) converts the received TF back to DD. The
transforms used in this project are
\begin{equation}
  x^{\mathrm{TF}}[n,m] \;=\;
  \frac{1}{\sqrt{N}}\sum_{l=0}^{N-1}
  \frac{1}{\sqrt{M}}\sum_{k=0}^{M-1}
  x^{\mathrm{DD}}[l,k]\,
  e^{\,j2\pi(nl/N - mk/M)}.
\end{equation}

\subsection{Embedded Pilot and Guard Region (TVT 2019)}
Following Raviteja \emph{et al.} (IEEE TVT 2019, Fig.~3), a single
impulse pilot of amplitude $x_p$ is placed at
$(l_p,k_p) = (l_\tau,\,M/2)$ in the DD grid. A guard region of rows
$l_p - l_\tau,\dots,l_p + l_\tau$ is reserved (no data), which
guarantees that the received pilot response and the data response do
not overlap. The pilot SNR is $\mathrm{SNR}_p = x_p^2/N_0$, set to
$45$\,dB in the simulation.

\subsection{Channel Estimation in the DD Domain}
After receiving the DD grid $Y_{\mathrm{DD}}$, the rows belonging to
the pilot response region (MATLAB rows $l_p+1,\dots,l_p+l_\tau+1$) are
divided by the pilot amplitude to obtain a noisy estimate of the
effective DD channel,
\begin{equation}
  \hat h_{\mathrm{DD}}[l,k]
  \;=\; \frac{Y_{\mathrm{DD}}[l,k]}{x_p},
  \qquad l\in\{l_p,\dots,l_p+l_\tau\}.
\end{equation}
A detection threshold
\begin{equation}
  \eta \;=\; \frac{3}{\sqrt{\mathrm{SNR}_p}}
\end{equation}
then zeros out bins dominated by noise. The surviving non-zero
entries form the sparse triplet set
$\{(\Delta l_s,\Delta k_s,\tilde h_s)\}$, where
$\Delta l_s = (l - l_p)\bmod N$ and
$\Delta k_s = (k - k_p)\bmod M$. These triplets are the effective
delay, Doppler and gain used by the detector.

\subsection{Message-Passing Demodulation}
Given the triplet channel, the detector of \texttt{mp\_detector.m}
implements the factor-graph message passing algorithm of Raviteja
\emph{et al.}\ (WCNC 2018, TVT 2019). For each data node $(l,k)$ with
candidate symbol $a_q$ from the 4-QAM alphabet, the variable-to-factor
messages carry the current mean $\mu[l,k]$ and variance $v[l,k]$;
factor-to-variable messages aggregate the interfering contributions of
the $S$ channel paths.

At each iteration the detector computes, for every observation bin,
the total expected signal and its variance,
\begin{equation}
  m_{\text{tot}}[l,k] = \sum_{s=1}^{S} \tilde h_s\,
  \mu\bigl[l-\Delta l_s,\,k-\Delta k_s\bigr],
  \qquad
  v_{\text{tot}}[l,k] = \sum_{s=1}^{S} |\tilde h_s|^2\,
  v\bigl[l-\Delta l_s,\,k-\Delta k_s\bigr],
\end{equation}
and then forms per-path \emph{extrinsic} observations by removing the
self-contribution,
\begin{equation}
  y_{\text{eff}}^{(s)}[l,k]
  = \frac{Y_{\mathrm{DD}}[l+\Delta l_s,k+\Delta k_s]
        - (m_{\text{tot}}[l+\Delta l_s,k+\Delta k_s] - \tilde h_s \mu[l,k])}
         {\tilde h_s},
\end{equation}
with effective noise variance
$\sigma^{2}_{\text{eff}}{}^{(s)}[l,k] = (N_0 + v_{\text{tot}} -
|\tilde h_s|^2 v)/|\tilde h_s|^2$. Symbol log-likelihoods are
accumulated over all paths,
\begin{equation}
  \log p_q[l,k] \;\propto\; -\sum_{s=1}^{S}
  \frac{|y_{\text{eff}}^{(s)}[l,k] - a_q|^{2}}
       {\sigma^{2}_{\text{eff}}{}^{(s)}[l,k]},
\end{equation}
converted to posteriors, damped with factor $\eta_d$ to stabilize
loopy propagation,
\begin{equation}
  p_q^{(t+1)}[l,k] \;=\; \eta_d\, p_q^{(t+1,\text{raw})}[l,k]
  \;+\; (1-\eta_d)\, p_q^{(t)}[l,k],
\end{equation}
and used to update $\mu[l,k]$ and $v[l,k]$ over the data mask only
(pilot and guard cells remain fixed). The simulation uses
$30$ iterations and $\eta_d = 0.7$. After convergence a hard decision
on $\log p_q$ produces $\hat X_{\mathrm{DD}}$, which is then QAM-
demodulated to bits.

\section{Maximum Likelihood Estimation of Path Parameters}

The MLE routine (\texttt{mle\_tf.m}) estimates the $3P$ unknowns
$\{\alpha_p,\tau_p,\nu_p\}_{p=1}^{P}$ from the observed TF grid
$Y_{\mathrm{TF}}$ and a (hard-decided) symbol hypothesis
$\hat X$. The ICI term is neglected
($\nu_p T \ll 1$), so the model used for estimation is
\begin{equation}
  Y[n,m] \;=\; \sum_{p=1}^{P}\alpha_p\,g_{n,m}(\tau_p,\nu_p) + Z[n,m],
  \qquad
  g_{n,m}(\tau,\nu) \;=\; \hat X[n,m]\,
  e^{-j2\pi \deltaf n \tau}\, e^{\,j2\pi T_s m \nu}.
  \label{eq:mle_model}
\end{equation}

Stacking columns, let $\mathbf y = \mathrm{vec}(Y)$ and
$\mathbf G(\bm\tau,\bm\nu) = [\mathbf g_1,\dots,\mathbf g_P]$ with
$\mathbf g_p = \mathrm{vec}(g_{:,:}(\tau_p,\nu_p))$. The negative log-
likelihood (up to additive constants) is
\begin{equation}
  \Lambda(\bm\alpha,\bm\tau,\bm\nu)
  \;=\; \bigl\|\mathbf y - \mathbf G(\bm\tau,\bm\nu)\bm\alpha\bigr\|^{2}.
\end{equation}
Since $\bm\alpha$ enters linearly, it is concentrated out (variable
projection) by the closed-form least-squares solution
\begin{equation}
  \hat{\bm\alpha}(\bm\tau,\bm\nu)
  \;=\; \mathbf G^{+}\mathbf y,
  \qquad
  \mathbf G^{+} = (\mathbf G^{\mathsf H}\mathbf G)^{-1}\mathbf G^{\mathsf H},
\end{equation}
reducing the cost to a $2P$-dimensional nonlinear problem in
$(\bm\tau,\bm\nu)$ only:
\begin{equation}
  \tilde\Lambda(\bm\tau,\bm\nu)
  \;=\; \bigl\|\bigl(\mathbf I - \mathbf G\mathbf G^{+}\bigr)\mathbf y\bigr\|^{2}.
\end{equation}

\subsection{Stage~1: RELAX/SIC Initialization}
A coarse-to-fine greedy search initializes the $P$ paths sequentially:
\begin{enumerate}
  \item Start with residual $\mathbf r = Y$.
  \item Form $A = \mathbf r \odot \hat X^{*}$.
  \item Compute a 2-D oversampled ambiguity surface
        $\mathrm{IFFT}_{N_f}(A)\cdot \mathrm{FFT}_{M_f}$ with
        oversampling factor $4$ in both axes. The peak gives
        $(\tau_p^{(0)},\nu_p^{(0)})$.
  \item Refine the peak with Nelder--Mead
        (\texttt{fminsearch}) on
        $|\sum_{n,m} A_{n,m}\,e^{+j2\pi\deltaf n\tau} e^{-j2\pi T_s m\nu}|^{2}$.
  \item Compute the closed-form complex gain
        $\hat\alpha_p = (\mathbf g_p^{\mathsf H}\mathbf r) /
        (\mathbf g_p^{\mathsf H}\mathbf g_p)$ and subtract
        $\hat\alpha_p \mathbf g_p$ from the residual.
\end{enumerate}

\subsection{Stage~2: Joint Refinement via Variable Projection}
The initial $2P$-vector $[\bm\tau^{(0)};\bm\nu^{(0)}]$ is then refined
by a joint nonlinear least-squares solver
(\texttt{lsqnonlin} with Levenberg--Marquardt). The residual fed to
the solver is
\begin{equation}
  \mathbf r(\bm\tau,\bm\nu)
  \;=\; \mathbf G(\bm\tau,\bm\nu)\,\hat{\bm\alpha}(\bm\tau,\bm\nu) - \mathbf y,
\end{equation}
evaluated with the concentrated $\hat{\bm\alpha}$. Real and imaginary
parts are stacked so that the LM Jacobian is well defined. At
convergence the path gains are recovered as
$\hat{\bm\alpha} = \mathbf G(\hat{\bm\tau},\hat{\bm\nu})\backslash \mathbf y$.
The two-stage structure mirrors the pseudo-true parameter solver used
by \texttt{getMCRB\_TF.m}, ensuring that the MLE and the MCRB
computation share a common model so that RMSE curves can be compared
directly against the mismatched CRB.

\subsection{CSI, CE, and Genie Variants}
Three variants of the MLE are evaluated in the BER/MCRB pipeline:
\begin{itemize}
  \item \textbf{Genie MLE}: run with the true transmitted symbols
        $X$ in place of $\hat X$. This is the matched CRB reference.
  \item \textbf{CSI MLE}: run with the hard-decided symbols from the
        MP detector operating on the true (thresholded) DD channel.
  \item \textbf{CE MLE}: run with the hard-decided symbols obtained
        from the detector that used only the channel estimated from
        the embedded pilot, i.e., a fully data-driven receiver.
\end{itemize}
Comparing the RMSE of these three MLEs against the genie CRB and the
MCRB (which accounts for the symbol mismatch) quantifies the loss
incurred by imperfect channel knowledge and by data-aided versus
pilot-aided processing.

\end{document}
